\section{Discussion}
\label{sec:discussion}

\subsection{What the Results Support}
\emph{Evaluation.} On these data the split protocol, not the model, decides whether scores look near-perfect; accuracy figures on gameplay cybersickness data should be read together with the split protocol.

\emph{RQ1.} On logged data the asymmetric reward encodes a cost-sensitive operating point: under the natural class mix the greedy policy is the 5:1 Bayes rule, and class-balanced replay shifts it. Empirically the ratio moved the default operating point but not the discrimination, and replacing the asymmetric reward with the symmetric one lowered rather than raised E-DDQN's cost as a point estimate (Section~\ref{sec:ablation}). The evidence therefore does not support a benefit of the asymmetric reward beyond what cost weighting and a validation-chosen threshold provide.

\emph{RQ2.} E-DDQN did not have a significantly lower 5:1 cost than the standard DQN with any of three search draws, and the two agents' ranking reversed between draws, so neither is consistently better. With its primary draw its seed-ensemble cost was higher than the random forest's, significantly under the primary procedure but not under the sensitivity analysis, so this is not firmly established. No primary model ranked seconds significantly better than elapsed time, and no model meets a safety-prioritized target.

\emph{RQ3.} The small effect of bounded perturbations on ranking reflects which inputs the selected configurations use; at fixed thresholds an adaptive adversary still changes decisions. These feature-space stress tests say nothing about real rendering attacks.

\subsection{Is E-DDQN Reinforcement Learning?}
With logged data the alarm cannot influence later telemetry, so the problem is a contextual bandit. E-DDQN remains a valid value-based learner of a cost-sensitive policy, but bootstrapping cannot improve the optimal policy and, under bandit feedback, the agent observes the reward of one action although the labels reveal both---which is why full-information replay is the principled way to use the labels. Any DQN classifier trained on a fixed dataset whose episodes end only with the data must therefore gain through estimation effects and should be compared with cost-weighted models at the same operating point; error-triggered termination~\cite{lin2020imbalanced,yang2024drlimbalanced} only adds a $\gamma$-dependent penalty on that error. Adaptive sequential control would require an alarm that triggers a mitigation whose effect on later telemetry and symptoms is measured, which, to our knowledge, no public cybersickness dataset provides.

\subsection{The Safety Objective}
A 5:1 cost makes an alarm cheaper whenever $P(y=1\mid s)\ge1/6$. With a prevalence of 0.44 and weakly informative inputs almost every second qualifies, so every cost-optimal policy is close to always alarming: high recall, but \DOneFPRMin--\DOneFPRMax\,\% false alarms. High recall alone is therefore not evidence of a useful monitor. Judged on cost, threshold-free ranking and cost curves across ratios, near-complete detection with a controlled false-alarm rate is not achieved on these data. Once sessions and groups are not shared, no primary model ranked seconds significantly better than session duration, which distinguishes individual susceptibility only coarsely; questionnaire indicators, which could in principle capture susceptibility, appeared not to transfer across the \nGroupsEval{} groups (exploratory; Section~\ref{sec:why}).

\subsection{Limitations}
\label{sec:limits}
(i)~One public dataset with \nSessionsEval{} sessions, \nGroupsEval{} groups and \nEventSessions{} onsets: intervals are wide, fold-to-fold variation is large, and inference with so few clusters is approximate (E-DDQN vs.\ random forest: Holm-adjusted $p=$\,\PHHolmRF{} by bootstrap, \PHSfRFHolm{} by sign-flip). (ii)~Participant identity cannot be verified: linkage groups are reconstructed from questionnaire profiles and session timing; the checks of Supplementary Sec.~S8 establish session- and group-disjoint evaluation, not participant independence. (iii)~The label is a coarse, held self-report with undocumented elicitation, and the binary threshold merges slight with severe discomfort, which is almost absent from Flight. (iv)~All analyses beyond the frozen protocol are post hoc; Section~\ref{sec:protocol} and Supplementary Sec.~S5 state which were committed before being run (the exposure-time baselines, Holm adjustment of the E-DDQN comparisons, always-alarm comparison, grouping sensitivity, cost curves and the search-draw experiment) and which were added later, some after a pre-submission review when all results were known. (v)~Perturbations are synthetic and feature-space, their physiological effect is unknown, and the adversary is weaker than optimization-based attacks. (vi)~The search budget (\HPOBudget{} configurations; 12 for logistic regression) covers a far smaller fraction of the RL spaces (about \PHSpaceEDDQNMillion{} million E-DDQN configurations) than of the random-forest space (\PHSpaceRF); only E-DDQN and the DQN were re-tuned with further draws, so the other rankings may also depend on their draws. (vii)~The October-2019 sessions could not serve as an external test set because their orientation channels are constant (a logging defect). (viii)~The diagnosis of E-DDQN's weak result rests on three draws and one feature-set restriction and does not isolate any single component. (ix)~Thresholds are chosen on OOF scores from single inner models but applied to test scores averaged over three inner models (and over seeds for the ensemble); averaging can shrink scores and contribute to imperfect threshold transfer. (x)~For models trained on a balanced class mix and for the standard DQN, the untuned rule D0 is the 5:1 Bayes cut-off only nominally; D1 and the ranking metrics are affected only marginally (Section~\ref{sec:formulation}).

\subsection{Implications}
Results on session data should be reported with participant- or session-grouped splits, trivial policies and an exposure-time baseline. RL detectors compared under small search budgets should be evaluated over several random draws; here E-DDQN's within-fold ROC-AUC spanned \SDWfSpan{} across three draws. An asymmetric reward should be treated as a cost-sensitive operating point and compared with cost-weighted models at matched recall or cost. Robustness claims should state which inputs a detector uses and evaluate decisions at fixed thresholds. Progress is likely to require physiological or eye-tracking signals, verified participant identifiers and, for adversarial claims, manipulations applied to participants with measured effects.
